\hypersetup{pdftitle={Submission 4: Part C Report},pdfauthor={}}
\begin{document}
\docheader{Submission 4: Part C, Experiment on a Real Online IR Dataset}
\section{Summary}
I ran TF-IDF, BM25 and a Dirichlet Language Model on the Cranfield test collection using Apache Lucene. With default parameters, BM25 (MAP [[c('BM25','AP')]]) and TF-IDF (MAP [[c('TF-IDF','AP')]]) were statistically indistinguishable, and both were significantly better than the Language Model (MAP [[c('LM-Dirichlet','AP')]]). Every value in this report comes from an actual run.

\section{Dataset description and source}
\begin{table}[H]\small\centering\caption{Cranfield collection.}
\begin{tabularx}{\linewidth}{@{}lY@{}}\toprule
Documents & [[D['docs']]] aeronautical engineering abstracts; [[IDX['tokens']]] tokens after analysis, [[IDX['vocab']]] unique terms, mean [[f"{float(IDX['meanLen']):.1f}"]] terms per document \\
Queries & [[D['queries']]] natural-language questions, mean [[f"{D['mean_query_words']:.1f}"]] words \\
Qrels & [[D['qrel_rows']]] rows, [[D['positive_rows']]] of them positive; mean [[f"{D['relevant_per_query_mean']:.2f}"]] relevant documents per query, from [[D['relevant_per_query_min']]] to [[D['relevant_per_query_max']]] \\
Why chosen & It is public, has the components the assignment requires, and is a historic benchmark [5]. TREC document sets usually require a license [6]. \\
Source used & TREC-format copy, commit 1208e6e: \url{https://github.com/oussbenk/cranfield-trec-dataset} \\
Original & \url{http://ir.dcs.gla.ac.uk/resources/test_collections/cran/} \\
\bottomrule\end{tabularx}\end{table}
I could not reach the original site from my working environment, so I used the TREC-format copy and checked it against the \texttt{ir\_datasets} metadata for Cranfield [7]. The document, query and qrel counts match (1400, 225, 1837), and the [[D['qrel_label_counts']['0']]] rows labeled 0 match the metadata's 225 ``no interest'' ($-1$) rows. The metadata lists grades 1 to 4 for the other rows, and this copy has flattened them to 1, apart from one row labeled 3 that I treated as relevant. Documents [[D['empty_text_docs'][0]]] and [[D['empty_text_docs'][1]]] have no abstract, so I indexed their titles.

\section{Experimental setup and methodology}
\begin{table}[H]\small\centering\caption{Setup, the same for all models.}
\begin{tabularx}{\linewidth}{@{}lY@{}}\toprule
Engine & Apache Lucene [[IDX['lucene']]], Java 21 (jar from Anserini 2.3.0 [8], distributed with Pyserini 2.4.0) \\
Analysis & \texttt{EnglishAnalyzer} on documents and queries (default stop words, Porter stemmer) \\
Indexed text & Abstract field, which already begins with the title in the source \\
Query & Topic text analyzed, then all terms combined with OR (\texttt{BooleanQuery}, SHOULD) \\
Depth & Top 1000; a mean of [[f"{C['retrieved_mean']['BM25']:.0f}"]] documents returned, since only documents with a query term match \\
Saved results & TREC-format run files (\texttt{results/partc\_*.run}) \\
Evaluation & Part B functions; set metrics on the full list and at rank 10; AP over the list; nDCG with binary gain \\
Testing & Paired two-sided randomization test (100,000 sign flips, seed 618) and Wilcoxon signed-rank test [9] \\
\bottomrule\end{tabularx}\end{table}
\begin{table}[H]\small\centering\caption{Models. $f$ is a query term's frequency in document $d$ and $|d|$ the document length.}
\renewcommand{\arraystretch}{1.3}
\begin{tabularx}{\linewidth}{@{}lYl@{}}\toprule
Model & Scoring idea & Parameters \\\midrule
TF-IDF (\texttt{ClassicSimilarity}) & $\sqrt{f}\cdot \mathrm{idf}^2\cdot |d|^{-1/2}$, summed over query terms & defaults \\
BM25 (\texttt{BM25Similarity}) [10] & idf times a saturating function of $f$, normalized by $|d|$ relative to the mean & $k_1$=1.2, $b$=0.75 \\
LM-Dirichlet (\texttt{LMDirichletSimilarity}) [11] & Term probability in $d$, smoothed with the collection model & $\mu$=2000 \\
\bottomrule\end{tabularx}\end{table}

\section{Results}
\begin{table}[H]\small\centering\caption{Mean metrics over [[D['queries']]] queries. The best value in each row is in bold; ties are not marked.}
\begin{tabular}{@{}lrrr@{}}\toprule
Metric & TF-IDF & BM25 & LM-Dirichlet \\\midrule
[[t_c()]]
\bottomrule\end{tabular}\end{table}
\begin{table}[H]\small\centering\caption{Paired tests over 225 queries. W/L/T counts queries where the first model scored higher, lower or equal.}
\begin{tabular}{@{}llrrrr@{}}\toprule
Comparison & Metric & Mean diff. & W/L/T & Randomization $p$ & Wilcoxon $p$ \\\midrule
[[t_tests()]]
\bottomrule\end{tabular}\end{table}
With a Bonferroni threshold of $0.05/12=0.0042$, all eight comparisons involving the Language Model stay significant and none of the BM25 vs TF-IDF comparisons is.

\begin{table}[H]\small\centering\caption{Query-level behavior. All three models tied for best AP on [[C['best_AP']['tie']]] queries.}
\begin{tabular}{@{}lrrr@{}}\toprule
 & TF-IDF & BM25 & LM-Dirichlet \\\midrule
Queries where the model alone has the best AP & [[C['best_AP']['TF-IDF']]] & [[C['best_AP']['BM25']]] & [[C['best_AP']['LM-Dirichlet']]] \\
Queries with no relevant document in the top 10 & [[C['zero_P10']['TF-IDF']]] & [[C['zero_P10']['BM25']]] & [[C['zero_P10']['LM-Dirichlet']]] \\
\bottomrule\end{tabular}\end{table}
Mean top-10 overlap was [[f"{C['overlap10']['BM25|TF-IDF']:.3f}"]] for BM25 and TF-IDF, [[f"{C['overlap10']['BM25|LM-Dirichlet']:.3f}"]] for BM25 and the Language Model, and [[f"{C['overlap10']['TF-IDF|LM-Dirichlet']:.3f}"]] for TF-IDF and the Language Model.

\begin{table}[H]\small\centering\caption{Five queries with the largest AP spread across models.}
\begin{tabularx}{\linewidth}{@{}rYrrrr@{}}\toprule
ID & Query & Rel. & TF-IDF & BM25 & LM \\\midrule
[[t_ex()]]
\bottomrule\end{tabularx}\end{table}

\begin{table}[H]\small\centering\caption{Parameter sensitivity. These settings were picked after viewing test queries, with no separate tuning set, so treat them as exploratory.}
\begin{tabular}{@{}lrrr@{}}\toprule
Configuration & MAP & MRR & nDCG@10 \\\midrule
[[t_sens()]]
\bottomrule\end{tabular}\end{table}

\section{Why performance differs}
BM25 and TF-IDF behave alike here. Their top-10 lists overlap by about three quarters, and both use IDF and length normalization. My guess is that Cranfield's short abstracts give BM25's term-frequency saturation few chances to matter, but I did not test this directly.

The Language Model is the outlier. MAP improves steadily as $\mu$ falls from 2000 to 50 ([[c('LM-Dirichlet','AP')]] to [[c('LMD-mu50','AP')]]), which fits over-smoothing, since $\mu$=2000 is about 20 times the mean document length. Tuning does not close the gap to BM25, and Jelinek-Mercer smoothing tops out at a similar value ([[c('LMJM-l0.7','AP')]]). Lucene's documentation says this class returns 0 where the formula would give a negative score, and only documents containing a query term are scored. Those choices may also limit it, but I did not isolate their effect.

The means hide differences between queries. On query 81, BM25's AP ([[f"{EX['81']['BM25']:.3f}"]]) is double TF-IDF's ([[f"{EX['81']['TF-IDF']:.3f}"]]), and on query 52 TF-IDF leads BM25 by [[f"{EX['52']['TF-IDF']-EX['52']['BM25']:.2f}"]].

\section{Verification}
My metrics matched the independent \texttt{ranx} library [12] to six decimal places on all eight metrics I compared (three are shown below). A hand check agrees too: for query 81 under BM25, relevant documents 799 and 672 are at ranks 1 and 3, so AP $=(1/1+2/3)/2=0.8333$ and RR $=1$, the same as the program.
\begin{table}[H]\small\centering\caption{My implementation vs ranx.}
\begin{tabular}{@{}lrrrrrr@{}}\toprule
Model & MAP (mine) & MAP (ranx) & MRR (mine) & MRR (ranx) & nDCG@10 (mine) & nDCG@10 (ranx) \\\midrule
[[t_ranx()]]
\bottomrule\end{tabular}\end{table}

\section{Limitations}
\begin{itemize}
\item The collection is small, old and domain-specific, so these rankings need not hold for other search tasks.
\item The qrels copy is binary, so nDCG here does not measure graded relevance. That would need the original graded file.
\item Unjudged documents count as non-relevant.
\item The main comparison uses default parameters, and the sensitivity table had no separate tuning set.
\item Lucene's versions differ from textbook forms, and the analyzer leaves question words such as ``what'' in queries.
\item There was no user study, and I did not compare speed or index size.
\end{itemize}

\section*{References}
{\small
[1] C. D. Manning, P. Raghavan, and H. Sch\"utze, \emph{Introduction to Information Retrieval}, Cambridge Univ. Press, 2008, Ch. 8.\par
[2] Apache Lucene, \url{https://lucene.apache.org/}.\par
[3] Cranfield, TREC-format copy, \url{https://github.com/oussbenk/cranfield-trec-dataset}.\par
[4] Cranfield, original files, \url{http://ir.dcs.gla.ac.uk/resources/test_collections/cran/}.\par
[5] C. Cleverdon, ``The Cranfield tests on index language devices,'' \emph{Aslib Proc.}, 19(6), 1967.\par
[6] NIST TREC FAQ, \url{https://trec.nist.gov/faq.html}.\par
[7] S. MacAvaney et al., ``Simplified data wrangling with ir\_datasets,'' SIGIR 2021.\par
[8] J. Lin et al., ``Pyserini: A Python toolkit for reproducible information retrieval research with sparse and dense representations,'' SIGIR 2021.\par
[9] M. Smucker, J. Allan, and B. Carterette, ``A comparison of statistical significance tests for information retrieval evaluation,'' CIKM 2007.\par
[10] S. Robertson and H. Zaragoza, ``The probabilistic relevance framework: BM25 and beyond,'' \emph{Found. Trends Inf. Retr.}, 3(4), 2009.\par
[11] C. Zhai and J. Lafferty, ``A study of smoothing methods for language models applied to information retrieval,'' \emph{ACM TOIS}, 22(2), 2004.\par
[12] E. Bassani, ``ranx: A blazing-fast Python library for ranking evaluation and comparison,'' ECIR 2022.\par
}

\clearpage\appendix\section*{Appendix A: Code}
The folder \texttt{part\_b\_c} in \url{https://github.com/Atmosu/csc618-ir-evaluation-rocchio-cranfield} holds this code, the main run files and the result tables.

\texttt{prepare\_cranfield.py}
\lstinputlisting[language=Python]{../code/prepare_cranfield.py}
\texttt{LuceneExperiment.java}
\lstinputlisting[language=Java]{../code/LuceneExperiment.java}
\texttt{ir\_metrics.py}
\lstinputlisting[language=Python]{../code/ir_metrics.py}
\texttt{part\_c\_eval.py}
\lstinputlisting[language=Python]{../code/part_c_eval.py}
\texttt{verify.py}
\lstinputlisting[language=Python]{../code/verify.py}
\end{document}
